% !PRINTED 127
\weberpage{127}

\begin{center}
{\large \S.\ 39.}

\vspace{0.5em}
{\large\gesperrt{Algorithmus zur Berechnung der Wurzeln.}}
\end{center}

Der im vorigen Paragraphen gegebene Beweis für die Existenz einer Wurzel einer algebraischen Gleichung lässt zwar an Bündigkeit nichts zu wünschen übrig, er hat aber noch den Mangel, dass er nicht die Schritte erkennen lässt, wie man eine Wurzel durch ein convergentes Rechnungsverfahren berechnen kann. Wir fügen also noch die folgenden Betrachtungen hinzu, die diesem Mangel, wenn auch nur theoretisch, abzuhelfen bestimmt sind. Sie geben noch einen zweiten Beweis des Fundamentalsatzes, der in der Hauptsache von \gesperrt{Lipschitz} herrührt (Lehrbuch der Analysis, Bd. I, \S.\,61 ff.; eine Vereinfachung verdanke ich einer Mittheilung von \gesperrt{Dedekind}).

Wenn die beiden ganzen rationalen Functionen $f(x)$ und $f'(x)$ einen gemeinsamen Theiler haben, so kann dieser nach \S.\,6 durch rationale Operationen gefunden und beseitigt werden. Wir dürfen also voraussetzen, dass $f(x)$ und $f'(x)$ keinen gemeinsamen Theiler haben, dass sie also für keinen Werth von $x$ beide zugleich verschwinden.

Wir setzen nun voraus, dass die Wurzelberechnung für eine Function $(n-1)$ten Grades schon gelungen sei, und nehmen demnach $f'(x)$ in lineare Factoren zerlegt an. Demnach sei, wenn
\begin{equation}
 f(x)=x^n+a_1x^{n-1}+a_2x^{n-2}+\cdots,
\tag{1}
\end{equation}
ist,
\begin{equation}
 f'(x)=nx^{n-1}+(n-1)a_1x^{n-2}+\cdots
      =n(x-\beta_1)(x-\beta_2)\cdots(x-\beta_{n-1}),
\tag{2}
\end{equation}
worin die $\beta_1,\beta_2,
\ldots,\beta_{n-1}$ als bekannte Zahlen angesehen werden, die auch theilweise identisch sein können. Nach unserer Voraussetzung werden die absoluten Werthe
\begin{equation}
 |f(\beta_1)|,\ |f(\beta_2)|,\ldots, |f(\beta_{n-1})|,
\tag{3}
\end{equation}
die wir mit
\begin{equation}
 b_1,b_2,\ldots,b_{n-1}
\tag{4}
\end{equation}
bezeichnen, alle von Null verschieden sein; wir wollen die Bezeichnung so gewählt annehmen, dass $b_1$ der kleinste unter ihnen
%CONT>>
